\documentclass[11pt]{article}
\usepackage{amsmath,amssymb,amsthm}
\usepackage[margin=1.2in]{geometry}
\usepackage{microtype}
\usepackage{parskip}

\begin{document}
\pagestyle{empty}

\begin{center}
  {\large\bfseries RetroDiff: Training Objective Derivation}
\end{center}

\medskip

We derive the training objective for RetroDiff by starting from the
standard variational lower bound on the log-likelihood of a diffusion
model.  Let $x$ denote a clean data sample, $x_t$ its noised version at
diffusion step $t \in \{1,\dots,T\}$, $q$ the fixed forward noising
process, and $p_\theta$ the learned reverse process with parameters
$\theta$.

\paragraph{Step 1 — Variational lower bound (ELBO).}
Applying Jensen's inequality to the intractable log-marginal
$\log p_\theta(x)$ yields the evidence lower bound:

\begin{align}
  \log p_\theta(x)
    &\;\geq\;
    \mathbb{E}_{q(x_{1:T}|x)}\!\left[
      \log\frac{p_\theta(x_{0:T})}{q(x_{1:T}|x)}
    \right]
    \;=\; -\mathcal{L}_{\mathrm{ELBO}}
  \label{eq:elbo}
\end{align}

Maximising the right-hand side is equivalent to minimising
$\mathcal{L}_{\mathrm{ELBO}}$, which we expand next.

\paragraph{Step 2 — Reconstruction + KL decomposition.}
Factoring the joint distributions and applying the Markov structure of
the forward process decomposes \eqref{eq:elbo} into a reconstruction
term and a sum of per-step KL divergences:

\begin{align}
  \mathcal{L}_{\mathrm{ELBO}}
    &= \underbrace{
         \mathbb{E}_q\!\left[-\log p_\theta(x_0 \mid x_1)\right]
       }_{\text{reconstruction}}
     + \sum_{t=2}^{T}
       \underbrace{
         \mathrm{KL}\!\left(
           q(x_{t-1}\mid x_t, x)
           \;\|\;
           p_\theta(x_{t-1}\mid x_t)
         \right)
       }_{\text{step-}t\text{ denoising}}
  \label{eq:decomp}
\end{align}

Each KL term measures how well the reverse step $p_\theta(x_{t-1}|x_t)$
matches the tractable posterior $q(x_{t-1}|x_t,x)$; the reconstruction
term handles the final generation step.

\paragraph{Step 3 — Retrieval-augmented objective (RetroDiff).}
RetroDiff's key contribution is to condition the entire ELBO on a set
$R(x)$ of $k$ exemplars retrieved from the training corpus based on
their similarity to $x$.  The per-sample objective is then averaged over
the data distribution:

\begin{align}
  \mathcal{L}_{\mathrm{RetroDiff}}(\theta)
    &= \mathbb{E}_{x \sim p_{\mathrm{data}}}\!\left[
         \mathcal{L}_{\mathrm{ELBO}}\!\left(x \;\big|\; R(x)\right)
       \right]
  \label{eq:retro}
\end{align}

Equation~\eqref{eq:retro} reduces to the standard diffusion objective
when $R(x) = \varnothing$, making RetroDiff a strict generalisation of
the baseline.  In practice $R(x)$ is computed once per sample with a
frozen retriever and cached for the duration of training.

\end{document}
